%% ML4PS 2026 submission -- 4 pages excluding references.
%% MUST use the NeurIPS 2026 LaTeX template with NO modifications
%% except the footer text below.
\documentclass{article}
\usepackage[final]{neurips_2026}
\usepackage[utf8]{inputenc}
\usepackage[T1]{fontenc}
\usepackage{hyperref}
\usepackage{url}
\usepackage{booktabs}
\usepackage{amsmath}
\usepackage{amssymb}
\usepackage{graphicx}
\usepackage{microtype}

% ---- REQUIRED footer override (the ONLY permitted modification) ----
\newcommand{\mlfpsfooter}{%
  \begin{center}\small
  Submitted to the 9th Workshop on Machine Learning and the Physical Sciences (ML4PS 2026). Do not distribute.
  \end{center}}

\title{What Really Limits Physics-Informed Operator Learning?\\[2mm]
\large A silent broadcasting pitfall, a training-free error-attribution diagnostic,\\
and evidence that subspace alignment regularizes rather than expands}

% Double-blind: no author information.
\author{}

\begin{document}
\maketitle
\mlfpsfooter

\begin{abstract}
Physics-informed operator learning is increasingly analysed through subspace diagnostics, with a
recurring claim that training fails because the trunk network's learned basis functions cannot span
the target solution space. We report three findings that redirect this line of inquiry.
\textbf{First}, a silent numerical pitfall: the routine that evaluates the input function at query
points is naturally written under a shared-grid convention returning shape \texttt{(n\_coef, *x.shape)};
when a batched loop hands it per-sample \texttt{(B, N)} coordinates, NumPy broadcasting silently
produces \texttt{(B, B, N)}, so each sample is trained against another sample's source term or
initial condition. It survives every conventional check---the dataset stays correct (ground-truth
finite-difference residual $6\!\times\!10^{-6}$), shapes are legal, supervised training succeeds.
Fixing it changes PI-DeepONet test error from \textbf{0.986} to \textbf{0.174} (diffusion--reaction)
and from \textbf{0.999} to \textbf{0.721} (Burgers). Critically, the buggy pipeline produced a clean,
monotone, publishable-looking narrative (Spearman $\rho=-1$ between subspace alignment and accuracy;
a dramatic ``trunk collapse'')---all artefactual.
\textbf{Second}, a training-free attribution diagnostic: with trunk features $\Phi$ and true
solutions $S$, the best relative error of \emph{any} branch network is exactly
$\sqrt{1-\text{captured\_energy}}$, so $\text{test\_relL2}=\mathrm{LB}+\mathrm{opt\_gap}$ separates
representation from optimization limits without training to convergence. After the fix
$\mathrm{LB}\approx0.19$--$0.25$ while $\mathrm{opt\_gap}\approx0.45$--$0.49$: representation is not
binding. \textbf{Third}, a mechanism: alignment of the trunk to an analytic spectral prior halves
loss sharpness and doubles output-NTK effective rank while captured energy moves only
$0.947\!\to\!0.951$; over \textbf{ten settings $\times$ five seeds} the gain is positive in $10/10$
and significant in $8/10$; and in 2D, where the baseline collapses deterministically onto $u\equiv0$
($\mathrm{relL2}=1.004$, NTK rank $1.1$) and is \textbf{representation-bound} ($85\%$ of its error is
the lower bound of Section~\ref{sec:attrib}), a random Fourier-feature trunk---no target subspace, no
tuning---repairs it better than alignment does ($1.004\!\to\!0.693$, \textbf{30.9\%}, versus $0.821$
for alignment, $p=1.1\times10^{-3}$), while a target-free decorrelation regularizer and a $10\times$
initial-condition weight both fail entirely. We also \textbf{retract} an earlier claim of ours
that the gain is linearly predicted by baseline instability: it collapses from $r=0.985$ to
$r=0.318$ (95\% CI $[-0.199,0.671]$) with better data.
\end{abstract}

\section{Introduction}
\label{sec:intro}

PI-DeepONet \cite{wang2021learning} and PINO \cite{li2021pino} train operator networks from governing
equations rather than paired simulation data. Because they are expensive and hard to debug, a
substantial literature explains their failures: spectral bias, gradient pathologies
\cite{wang2021understanding}, causality violations \cite{wang2022respecting}, landscape stiffness
\cite{krishnapriyan2021characterizing}, and---most relevant here---the proposal that the trunk basis
cannot span the target solution space \cite{jcp2026deeponets}.

We began from that last hypothesis, built the full pipeline (a subspace diagnostic, a differentiable
alignment loss toward an analytic spectral prior, and controlled comparisons), and obtained excellent
results: trunk effective rank $1.5\!\to\!33.7$, captured energy $0.976\!\to\!0.9999$, mean principal
angle $67^\circ\!\to\!17^\circ$, with five methods ranked exactly as theory predicted. All of it was
an artefact of a one-line shape bug in our own PDE interface.

We therefore write the paper in two halves. Section~\ref{sec:pitfall} documents the pitfall
precisely, shows why standard sanity checks are blind to it, and demonstrates that it can
manufacture a coherent publishable story. Sections~\ref{sec:attrib}--\ref{sec:2d} report what
survives after the fix: an attribution diagnostic, and a mechanistic account of subspace alignment
as regularization. We regard the first half as the more immediately useful contribution to the
community.

\paragraph{Setup.} Burgers $u_t+uu_x=\nu u_{xx}$ on $x\in[0,1)$ periodic, $t\in[0,1]$, solved by a
Strang-split spectral method; also 2D Burgers ($u_t+u\cdot\nabla u=\nu\Delta u$, 2D spectral
splitting) and a 1D diffusion--reaction system. Input functions are truncated spectral expansions,
so the input field is analytic and available at any query point. Branch/trunk MLPs, 3 hidden layers,
Adam $\mathrm{lr}=10^{-3}$, 6000 steps (10000 in Sec.~\ref{ssec:artefact}), 64 collocation points
$\times$ 16 input functions per step, 200/50/50 train/val/test input functions. We report test
relative $L_2$ and use Welch's $t$-test over seeds.

\section{A silent pitfall in DeepONet-style training loops}
\label{sec:pitfall}

\subsection{The code pattern}
\label{ssec:pattern}

An operator-learning problem must evaluate the input function at arbitrary query points---a source
term in the residual, an initial-condition target, or a spatially varying coefficient. A natural
helper is

\begin{quote}
\begin{verbatim}
def input_at(self, coef, x):
    """coef: (n, K) coefficients of the input function.  x: query points."""
    S = np.sin(np.outer(x.ravel(), self.k))      # shared-grid convention
    C = np.cos(np.outer(x.ravel(), self.k))
    out = coef @ S.T + coef @ C.T                # (n, M)
    return out.reshape(coef.shape[0], *x.shape)
\end{verbatim}
\end{quote}

This is correct whenever $\texttt{x}$ describes one shared grid---sensor locations $(m,)$, the
evaluation grid $(N_g,)$, or a sensor mesh $(m,2)$. It is \textbf{wrong} when $\texttt{x}$ is a
per-sample array $(B,N)$, which is exactly what a batched loop hands over. With
$\texttt{coef.shape[0]}=\texttt{x.shape[0]}=B$ the output becomes $(B,B,N)$; every downstream
expression broadcasts, the loss stays a scalar of plausible magnitude, and sample $b$ is trained
against sample $b'$'s data for $B-1$ of $B$ pairings. There is no exception, no warning, and no NaN.
We found this in three independent PDEs: the Burgers initial-condition target, the
diffusion--reaction collocation source, and the Darcy coefficient together with its gradients.

\subsection{Why the usual checks are blind}
\label{ssec:blind}

All four checks a practitioner would run pass. (i)~\emph{Ground-truth residual}: finite-differencing
the stored solutions and evaluating the PDE gives $\mathrm{mean}(\mathrm{res}^2)=6.0\times10^{-6}$---
because the \emph{dataset is correct}; the reference solvers consume coefficients with per-sample
semantics and never call the buggy helper. (ii)~\emph{Supervised training} succeeds, reaching
$\mathrm{relL2}=0.26$, because data-driven training never calls the helper on per-sample coordinates.
(iii)~\emph{Shape assertions} pass: $(B,B,N)$ is a legal shape and broadcasting is silent.
(iv)~\emph{The loss magnitude} is plausible, since the mean over $B^2$ pairings is statistically
similar to the correct mean.

The single diagnostic that exposed the bug was the \textbf{autograd PDE residual of a supervised
model}: after fitting the data to $\mathrm{relL2}=0.26$, the model's own residual was
$\mathrm{mean}(\mathrm{res}^2)=0.368$, while the trivial solution $s\equiv0$ has residual
$\mathbb{E}[u^2]=0.185$. A model that reproduces the solution cannot have a residual worse than the
zero function; that contradiction was the thread that led to the bug.

\paragraph{Detection recipe.} Two lines, both cheap, which we recommend as standard practice:
\begin{quote}
\begin{verbatim}
assert res_autograd(model) < 0.5 * res_trivial   # after supervised training
assert pts["u0"].shape == (B, n_col)             # never (B, B, n_col)
\end{verbatim}
\end{quote}

\subsection{What the bug does to scientific conclusions}
\label{ssec:artefact}

This is the part we consider most important. Before the fix, our Burgers ablation (width 64, 10\,000
steps) looked like Table~\ref{tab:artefact}. Accuracy is monotone in captured energy ($\rho=-1$
across five methods); spectral priors beat a random subspace and a target-free regularizer; a clear
``trunk basis collapse'' (effective rank $1.16$) appears to explain the baseline's failure; and the
story ``PI-DeepONet fails because its trunk basis is misaligned'' looks supported.
\textbf{Every one of these conclusions is false.}

After fixing the one-line broadcasting bug---with architecture, data, optimizer and step count all
held fixed---test relL2 falls from $0.9857$ to $\mathbf{0.1737}$ on diffusion--reaction and from
$0.9992$ to $\mathbf{0.7213}$ on Burgers, while trunk captured energy rises from $0.9774$ to
$0.9999$ and from $0.6653$ to $0.9471$ respectively. The ``collapse'' essentially disappears: the
plain baseline reaches captured energy $0.947$ with no representation-side intervention at all.
Subspace-level conclusions about operator learning must not be drawn before the physics residual is
verified to actually be minimised.

\begin{table}[t]
\caption{\textbf{Artefactual} results produced by the buggy pipeline (Burgers, width 64, 10\,000
steps). All qualitative readings of this table are false.}
\label{tab:artefact}
\vskip 0.1in
\begin{center}
\begin{small}
\begin{tabular}{lcccccc}
\toprule
method & test relL2 & LB & opt\_gap & captured & eff.\ rank & angle \\
\midrule
mlp baseline         & 0.9992 & 0.5786 & 0.421 & 0.6653 & 1.16  & 74.4$^\circ$ \\
orthogonalization    & 0.9966 & 0.4826 & 0.514 & 0.7671 & 5.05  & 68.8$^\circ$ \\
random-subspace      & 0.9826 & 0.3366 & 0.646 & 0.8867 & 13.23 & 56.3$^\circ$ \\
spectral-prior       & 0.9809 & 0.2960 & 0.685 & 0.9124 & 6.18  & 52.5$^\circ$ \\
Fourier-feature trunk& 0.9493 & 0.2838 & 0.666 & 0.9195 & 4.81  & 50.8$^\circ$ \\
\bottomrule
\end{tabular}
\end{small}
\end{center}
\vskip -0.1in
\end{table}

\section{An error-attribution diagnostic}
\label{sec:attrib}

Let $\Phi=\mathrm{trunk}(Y)\in\mathbb{R}^{N\times p}$ be trunk features on a probe set and
$S\in\mathbb{R}^{N\times n}$ the matrix of true solutions.

\begin{proposition}
\label{prop:lb}
$\displaystyle \min_B \frac{\|\Phi B-S\|_F}{\|S\|_F}
 = \frac{\|(I-P_\Phi)S\|_F}{\|S\|_F}
 = \sqrt{1-\mathrm{captured\_energy}}$,
where $P_\Phi$ projects onto $\mathrm{span}(\Phi)$ and
$\mathrm{captured\_energy}=\|P_\Phi S\|_F^2/\|S\|_F^2$.
\end{proposition}
\begin{proof}
The minimiser is $B=\Phi^{+}S$, giving the projection residual; $\|P_\Phi S\|_F^2/\|S\|_F^2$ is
exactly the captured energy (energy-weighted by $\sigma_i^2$ in the version used throughout).
\end{proof}

Hence $\mathrm{LB}:=\sqrt{1-\mathrm{captured\_energy}}$ is a \emph{hard} lower bound on the error of
any branch network on this trunk, computable at any point during training from one SVD, and
$\mathrm{opt\_gap}:=\mathrm{test\_relL2}-\mathrm{LB}$ isolates the optimization gap. Large
$\mathrm{LB}$ $\Rightarrow$ representation-limited; large $\mathrm{opt\_gap}$ $\Rightarrow$
optimization-limited.

After the fix, all representation-side methods land at $\mathrm{captured}\in[0.94,0.96]$, i.e.\
$\mathrm{LB}\in[0.19,0.25]$, while $\mathrm{opt\_gap}\in[0.45,0.49]$ dominates---correctly
predicting that representation-side work has little headroom. Consistently, the two standard
optimization-side remedies both \emph{fail} here. \textbf{Causal weighting}
\cite{wang2022respecting} degrades monotonically ($0.721\to0.760\to0.805\to0.829$ for
$\varepsilon=1,10,100$): with a soft initial-condition constraint of equal weight, concentrating
residual weight at small $t$ starves the later-time fit. \textbf{L-BFGS fine-tuning} after Adam,
standard in single-instance PINN solvers, also degrades ($0.721\to0.912$; with a smaller fixed
training set, $\to2.004$): L-BFGS requires a deterministic objective over a fixed training set that
is far too small to represent an operator family (hundreds of input functions $\times$ thousands of
coordinates), and the second-order method overfits it.

\section{Mechanism: alignment regularizes rather than expands}
\label{sec:mechanism}

We measure two quantities that are \textbf{never optimized} and independent of the diagnostic used to
build the alignment loss: the output neural tangent kernel $K=JJ^\top$ (per-point Jacobian of the
model output w.r.t.\ all parameters) and loss sharpness $\lambda_{\max}(\nabla^2 L)/2L$ (power
iteration with Hessian-vector products) \cite{jacot2018ntk,keskar2017large}.

\begin{table}[t]
\caption{Paired comparison at width 64 (same seed, data, optimizer). Alignment halves sharpness and
doubles output-NTK effective rank, while captured energy barely moves.}
\label{tab:ntk}
\vskip 0.1in
\begin{center}
\begin{small}
\begin{tabular}{lccccc}
\toprule
checkpoint & trunk eff.\ rank & out-NTK eff.\ rank & NTK top-eig.\ share & $\lambda_{\max}$ & sharpness \\
\midrule
baseline, seed 1       & 4.1  & 1.7 & 0.874 & 1035 & 295 \\
\textbf{aligned}, seed 1 & \textbf{9.0} & \textbf{3.8} & \textbf{0.542} & \textbf{430} & \textbf{141} \\
baseline, seed 2       & 4.2  & 2.0 & 0.803 & 2020 & 682 \\
\textbf{aligned}, seed 2 & \textbf{10.6} & \textbf{4.0} & \textbf{0.432} & \textbf{833} & \textbf{336} \\
\bottomrule
\end{tabular}
\end{small}
\end{center}
\vskip -0.1in
\end{table}

Alignment roughly \textbf{doubles the output-NTK effective rank}, cuts the dominant eigenvalue share
from $\approx0.85$ to $\approx0.49$, and \textbf{halves sharpness}, with $\lambda_{\max}$ falling from
959--2020 (all baselines $\ge959$) to 430--833 (both aligned $<959$)---a complete separation between
conditions. Captured energy moves only $0.947\!\to\!0.951$: this is \textbf{not} a capacity
expansion.

\subsection{Ten settings, five seeds}
\label{ssec:phase}

\begin{table}[t]
\caption{Ten settings, five seeds each (6000 steps). The gain is positive in $10/10$ and significant
in $8/10$. Note the two-sided variance clamp in the last column's $\sigma$ values.}
\label{tab:phase}
\vskip 0.1in
\begin{center}
\begin{small}
\begin{tabular}{llccc}
\toprule
setting & baseline (mean $\pm$ \textbf{std}) & aligned (mean $\pm$ std) & gain & Welch $p$ \\
\midrule
\textbf{2D, w32}   & 1.0039 $\pm$ 0.0024 & 0.8205 $\pm$ 0.0421 & \textbf{18.27\%} & 0.0006 \\
$\nu{=}0.02$, w32  & 0.8047 $\pm$ \textbf{0.0780} & 0.6815 $\pm$ 0.0102 & 15.31\% & 0.024 \\
$\nu{=}0.01$, w32  & 0.8353 $\pm$ \textbf{0.0729} & 0.7420 $\pm$ 0.0080 & 11.18\% & 0.045 \\
$\nu{=}0.02$, w64  & 0.7627 $\pm$ 0.0122 & 0.6823 $\pm$ 0.0118 & 10.55\% & $<$0.0001 \\
$\nu{=}0.005$, w32 & 0.8435 $\pm$ \textbf{0.0731} & 0.7683 $\pm$ 0.0056 & 8.92\%  & 0.083 \\
$\nu{=}0.02$, w24  & 0.7500 $\pm$ 0.0144 & 0.6874 $\pm$ 0.0081 & 8.35\%  & 0.0001 \\
$\nu{=}0.01$, w64  & 0.7999 $\pm$ 0.0033 & 0.7465 $\pm$ 0.0100 & 6.67\%  & 0.0001 \\
$\nu{=}0.005$, w64 & 0.8117 $\pm$ 0.0089 & 0.7691 $\pm$ 0.0091 & 5.25\%  & 0.0001 \\
$\nu{=}0.02$, w16  & 0.7370 $\pm$ 0.0145 & 0.7020 $\pm$ 0.0108 & 4.74\%  & 0.003 \\
$\nu{=}0.005$, w16 & 0.7974 $\pm$ 0.0128 & 0.7797 $\pm$ 0.0138 & 2.22\%  & 0.070 \\
\bottomrule
\end{tabular}
\end{small}
\end{center}
\vskip -0.1in
\end{table}

Table~\ref{tab:phase} shows the gain is \textbf{positive in $10/10$} settings and \textbf{significant
in $8/10$} (five at $p<0.001$). Its most robust effect is on variance: across the nine 1D settings the
baseline $\sigma$ spans \textbf{23.6$\times$} ($0.0033$--$0.0780$) while the aligned $\sigma$ spans
only \textbf{2.5$\times$} ($0.0056$--$0.0138$). Alignment therefore does not reduce variance by a
fixed proportion; it \textbf{clamps $\sigma$ into a problem-independent band
$\approx0.006$--$0.014$, from both sides}---$7.6$--$13.1\times$ reductions where the baseline is far
above it, $\approx1$ where the baseline already lies inside it, and $0.33\times$ where the baseline
(at $0.0033$) lies \emph{below} it, though the mean still improves by $6.67\%$ there. This two-sided
clamp holds in every 1D cell and fits better than any monotone law.

\paragraph{Retraction.} In an earlier iteration of this project we reported that the gain was almost
perfectly predicted by baseline $\sigma$ ($r=0.985$, four settings, three seeds). That claim does not
survive better data: with ten settings and five seeds it gives \textbf{$r=0.318$, 95\% seed-level
bootstrap CI $[-0.199,0.671]$}, stable as $n$ grew $8\to10$. The regressor is itself estimated from
few seeds, and fitting a noisy regressor inflates $|r|$. \textbf{We retract the linear law.} The
defensible claims are the sign of the gain, its significance, the variance clamp, and
Section~\ref{sec:2d}.

\section{2D: representation-bound, but alignment is not the fix}
\label{sec:2d}

In 2D the baseline does not partially fail---it \textbf{collapses completely and reproducibly}:
$1.0039\pm0.0024$ over five seeds (captured $\approx0.27$, out-NTK effective rank $1.08$--$1.10$, top
eigenvalue carrying $98.5$--$98.9\%$ of the spectrum). The model has collapsed onto $u\equiv0$, an
exact solution of the PDE and the periodic boundary conditions, leaving only the initial condition to
supply gradient; it is not \emph{unstable}---$\sigma=0.0024$---it is \emph{degenerate}. Splitting the
error three ways (Section~\ref{sec:attrib}) shows the failure is \textbf{representational}:
$\mathrm{LB}=0.849$ accounts for $85\%$ of the baseline's error while $\mathrm{opt\_gap\_train}=0.154$
and $\mathrm{gen\_gap}\approx0$.

We then isolated what repairs it (width 32, 6000 steps, five seeds):

\begin{table}[h]
\centering
\small
\begin{tabular}{lcccc}
\toprule
intervention & test relL2 & gain & $p$ & captured \\
\midrule
baseline & $1.0039\pm0.0024$ & --- & --- & 0.28 \\
orthogonalization reg.\ (target-free) & $1.0016\pm0.0010$ & 0.22\% & 0.112 & 0.30 \\
$10\times$ initial-condition weight & $1.0029\pm0.0013$ & 0.10\% & 0.464 & 0.28 \\
alignment, random smooth subspace & $0.7910\pm0.0165$ & 21.21\% & $6.0\!\times\!10^{-6}$ & 0.60 \\
alignment, analytic spectral prior & $0.8205\pm0.0421$ & 18.27\% & $6.1\!\times\!10^{-4}$ & 0.61 \\
\textbf{random Fourier-feature trunk} & $\mathbf{0.6933\pm0.0199}$ & \textbf{30.94\%} & $3.1\!\times\!10^{-6}$ & \textbf{0.77} \\
\bottomrule
\end{tabular}
\end{table}

\textbf{(i)} Diversity alone is not enough: the target-free decorrelation regularizer (only
$\Phi^\top\Phi\propto I$) is ineffective here ($p=0.112$)---what matters is which function subspace
the trunk spans, not how many dimensions. \textbf{(ii)} The prior does not matter: a random smooth
subspace matches the analytic spectral prior (0.791 vs 0.821). \textbf{(iii)} The best fix is free:
Fourier features beat alignment by 15.5\% ($p=1.1\times10^{-3}$). 2D is representation-bound, and what
rescues it is a \emph{spectrally structured trunk}, which architecture supplies more cheaply than an
alignment loss.

\section{Conclusion}
\label{sec:conclusion}

We set out to test whether PI-DeepONet fails because its trunk basis cannot span the solution space,
obtained a clean and convincing affirmative answer, and discovered that the answer was produced by a
one-line NumPy broadcasting bug that survives every conventional sanity check---including the
finite-difference residual of the ground truth, which stays at $6\times10^{-6}$ while training is
entirely broken. After the fix the representation-side story collapses: $\mathrm{LB}$ is already
small, the trunk is fine, and the residual error is optimization-bound---and the two standard
optimization-side remedies both make things worse. What survives is a training-free attribution
diagnostic that makes the representation/optimization split quantitative, and a mechanism---alignment
regularizes, and in 2D is necessary. Along the way we also retracted one of our own quantitative
claims, which is the same failure mode as the bug: a convincing regularity produced by too little
data and a noisy regressor, caught only because we went back and measured it properly. We hope the
detection recipe and the attribution diagnostic are useful beyond this paper's own claims.

\paragraph{Reproducibility.} Code, cached datasets, per-run JSON and all figures accompany this
submission as an anonymized release. Every number in every table is produced by the released scripts
(\texttt{scripts/run.py}, \texttt{scripts/analyze\_phase.py} with seed-level bootstrap,
\texttt{scripts/ntk\_analysis.py}); the per-run JSON files are included so that all tables and
figures can be regenerated from raw results without retraining.

\paragraph{Use of generative AI.} AI assistance was used for code drafting and for editing prose.
All numerical results were produced by executing the released scripts on the hardware described, and
all reported tables and statistics were recomputed from the accompanying per-run JSON files rather
than transcribed from generated text. Two specific correctness checks were performed on the AI-assisted
material: (i) the PDE residual of a supervised model was verified numerically against a trivial
solution baseline (this is how the pitfall in Section~\ref{sec:pitfall} was found), and (ii) the
statistical claims were independently recomputed with a corrected implementation after an error was
detected, leading to the retraction in Section~\ref{ssec:phase}.

\bibliographystyle{unsrt}
\begin{thebibliography}{16}
\bibitem{raissi2019pinn} M.~Raissi, P.~Perdikaris, G.~E.~Karniadakis. Physics-informed neural networks.
  \emph{J. Comput. Phys.}, 378:686--707, 2019.
\bibitem{lu2021deeponet} L.~Lu, P.~Jin, G.~Pang, Z.~Zhang, G.~E.~Karniadakis. Learning nonlinear
  operators via DeepONet. \emph{Nature Machine Intelligence}, 3(3):218--229, 2021.
\bibitem{wang2021learning} S.~Wang, H.~Wang, P.~Perdikaris. Learning the solution operator of
  parametric PDEs with physics-informed DeepONets. \emph{Science Advances}, 7(40):eabi8605, 2021.
\bibitem{li2021pino} Z.~Li, N.~Kovachki, K.~Azizzadenesheli, et~al. Physics-informed neural operator
  for learning PDEs. arXiv:2111.03794; \emph{ACM TOMS}, 2024.
\bibitem{wang2022respecting} S.~Wang, S.~Sankaran, P.~Perdikaris. Respecting causality when training
  physics-informed neural networks. arXiv:2203.07404, 2022.
\bibitem{krishnapriyan2021characterizing} A.~Krishnapriyan, A.~Gholami, S.~Zhe, R.~Kirby,
  M.~W.~Mahoney. Characterizing possible failure modes in physics-informed neural networks.
  \emph{NeurIPS}, 2021.
\bibitem{jcp2026deeponets} What do physics-informed DeepONets learn? Understanding and improving
  training. \emph{J. Comput. Phys.}, 2026. (S0021999126002019)
\bibitem{chen2026pino} N.~Chen, C.~Cui, A.~Chen, S.~Wang, R.~Ma. On the training of
  physics-informed neural operators for parametric PDEs. arXiv:2606.06164, 2026.
\bibitem{wang2022ntk} S.~Wang, X.~Yu, P.~Perdikaris. When and why PINNs fail to train: A neural
  tangent kernel perspective. \emph{J. Comput. Phys.}, 449:110768, 2022.
\bibitem{wang2021understanding} S.~Wang, Y.~Teng, P.~Perdikaris. Understanding and mitigating
  gradient pathologies in physics-informed neural networks. \emph{SIAM J. Sci. Comput.},
  43(5):A3055--A3081, 2021.
\bibitem{jacot2018ntk} A.~Jacot, F.~Gabriel, C.~Hongler. Neural tangent kernel. \emph{NeurIPS}, 2018.
\bibitem{keskar2017large} N.~S.~Keskar, D.~Mudigere, J.~Nocedal, et~al. On large-batch training:
  Generalization gap and sharp minima. \emph{ICLR}, 2017.
\bibitem{tancik2020fourier} M.~Tancik, P.~P.~Srinivasan, B.~Mildenhall, et~al. Fourier features let
  networks learn high frequency functions. \emph{NeurIPS}, 2020.
\bibitem{hu2026realpdebench} P.~Hu, H.~Feng, H.~Liu, et~al. RealPDEBench. \emph{ICLR}, 2026 (Oral);
  arXiv:2601.01829.
\bibitem{zhang2026pinns} W.~Zhang, W.~Suo, J.~Song, W.~Cao. PINNs as intelligent computing technique
  for PDEs: Limitation and future prospects. \emph{Sci. China Phys. Mech. Astron.}, 69:214602, 2026.
\bibitem{lu2021deepxde} L.~Lu, X.~Meng, Z.~Mao, G.~E.~Karniadakis. DeepXDE. \emph{SIAM Review},
  63(1):208--228, 2021.
\end{thebibliography}

\end{document}
